% Generated draft: requires mathematical and source review.
\documentclass[11pt]{article}
\usepackage[margin=25mm]{geometry}
\usepackage{fontspec}
\setmainfont{Latin Modern Roman}
\usepackage{amsmath,amssymb,mathtools,unicode-math}
\setmathfont{Latin Modern Math}
\usepackage{enumitem,needspace,xurl,hyperref}
\setlength{\parindent}{0pt}
\setlength{\parskip}{5pt}
\setlength{\emergencystretch}{3em}
\newcommand{\N}{\mathbb{N}}
\newcommand{\Z}{\mathbb{Z}}
\newcommand{\Q}{\mathbb{Q}}
\newcommand{\R}{\mathbb{R}}
\newcommand{\C}{\mathbb{C}}
\newcommand{\F}{\mathbb{F}}
\newcommand{\A}{\mathbb{A}}
\newcommand{\PP}{\mathbb{P}}
\newcommand{\E}{\mathbb{E}}
\newcommand{\eps}{\varepsilon}
\newcommand{\dd}{\,\mathrm{d}}
\begin{document}
\section*{UnsolvedMath ID: 30005071}
\textbf{Source problem code:} catalogue entry
\subsection{Definitions and mathematical statement}
\textbf{Formulation limitation.} The supplied record does not determine a precise mathematical conjecture. Its reconstructed statement describes a research programme for affine vertex algebras at admissible level, but does not specify the underlying Lie algebra, the admissible levels under consideration, the choice of universal or simple affine vertex algebra, or the category of modules. The retained original passage merely announces a discussion of progress, difficulties, and open questions. Consequently, the following gives a conditional formulation of the requested objectives, rather than a conjecture quantified over all affine vertex algebras or all their modules.

Let \(V\) denote an affine vertex algebra at a level called admissible in the supplied record, and let \(\mathcal C\) denote an intended abelian category of \(V\)-modules. Both \(V\) and \(\mathcal C\), including the meaning of admissibility and the restrictions on objects and morphisms, remain to be specified from a reliable mathematical source.

An object \(S\in\mathcal C\) is simple if it is nonzero and has no nonzero proper subobject. An object \(P\) is projective if every morphism from \(P\) to the target of an epimorphism lifts through that epimorphism. A projective cover of \(S\) is an epimorphism \(p:P\twoheadrightarrow S\), with \(P\) projective, such that any subobject \(U\subseteq P\) satisfying \(p(U)=S\) equals \(P\). Dually, an object \(I\) is injective if every morphism into \(I\) from the domain of a monomorphism extends across that monomorphism. An injective hull of \(S\) is a monomorphism \(S\hookrightarrow I\), with \(I\) injective, whose image meets every nonzero subobject of \(I\) nontrivially. Covers and hulls depend on the chosen category; their existence is an objective, not an assumption established by the record.

An object \(M\) has finite length if it admits a finite filtration
\[
0=M_0\subsetneq M_1\subsetneq\cdots\subsetneq M_n=M
\]
with simple successive quotients, allowing \(n=0\) when \(M=0\). A rigid braided tensor category has a monoidal product, unit, associativity and unit constraints, a braiding satisfying the hexagon identities, and left and right duals for every object satisfying the duality identities.

\textbf{Conditional research objectives.} Once the missing choices and hypotheses are fixed, the reconstructed statement asks to classify the simple objects of \(\mathcal C\) up to isomorphism; establish the existence and determine the projective covers and injective hulls of those simple objects; establish that every projective object of \(\mathcal C\) has finite length; and construct a rigid braided tensor-category structure on the intended category together with explicit fusion rules. The source does not specify whether all these objectives concern the same category, nor whether its phrase about projectives concerns every projective object or only the projective covers under discussion.

The meaning of fusion rules also requires specification: decomposition into direct summands, composition-factor multiplicities, and multiplication in a Grothendieck ring are different possible targets. No choice among them, no tensor-product construction, and no universal quantifier over Lie algebras or admissible levels can be supplied from this record alone.
\subsection{Short English statement}
The reconstructed problem asks for a description of simple modules, their projective covers and injective hulls, finite-length results for projectives, and rigid braided tensor structures with computable fusion rules for affine vertex algebras at admissible level. The vertex algebras, module categories, and scope of these requests are unspecified, so this remains a conditional research programme requiring source clarification and human mathematical review.
\subsection{Sources}
\begin{itemize}
\item Local UnsolvedMath registry, supplied complete source record: permanent ID \(30005071\), title \emph{Module Categories of Affine Vertex Algebras}, source problem code \texttt{catalogue entry}, collection \emph{Oberwolfach Reports Open Problems}. No individual proposer is supplied; the proposed year is recorded as 2022.
\item Supplied report citation: \emph{Mini-Workshop: Recent Developments in Representation Theory and Mathematical Physics (2023)}, Oberwolfach Reports, pp.~817--848, DOI \texttt{10.4171/owr/2022/16}. Supplied source URL and literature-source link, labelled \emph{Oberwolfach Report 2022/16}: \url{https://doi.org/10.4171/owr/2022/16}. The citation's parenthetical year 2023 and the report identifier's year 2022 are retained without reconciliation. No passage locator for the reconstructed mathematical objectives is supplied.
\item The retained original statement announces a summary of recent progress, difficulties, and open questions, and describes the report passage as retained context. The background identifies Oberwolfach Reports as reports of mathematical workshops and their open-problem discussions. The record separately labels the reconstruction \texttt{corrected\_verified} and states that it was checked against the official report; these are supplied provenance claims, not verification performed for this draft.
\item Literature assessment dated exactly \texttt{2026-08-22}: the supplied status is \emph{open}, while the supplied assessment says that the extracted fragment is meeting prose rather than a mathematically identifiable claim that can be verified or refuted. This is supplied dated curation derived from the supplied OWR curation, not independently verified. The record requests independent verification of time-sensitive status and cited claims; none is asserted here.
\end{itemize}
\end{document}
